\section{Introduction}

Vision--language models (VLMs) can answer questions about objects, visible text, and spatial relations \citep{antol2015vqa,okvqa,schwenk2022aokvqa}. Yet a correct answer does not show whether the model used the image region that contains it: a typical object color or the surrounding scene can suggest the same response. Language priors and the limits of saliency-based explanations make this uncertainty consequential \citep{goyal2017vqav2,agrawal2018vqacp,adebayo2018sanity,jain2019attention}. Recent VLM studies also locate consequential visual information and trace interpretable internal structures \citep{golovanevsky2025notice,li2026causalTracing,yang2026vlmcircuittracing}. These observations occupy different evidentiary levels. A region may be associated with an answer; changing it may alter the answer; and a proposed internal object may or may not reproduce the score change caused by that alteration. We ask: \emph{when compact answer-bearing evidence is removed, which internal descriptions account for the resulting answer-score change?}

The first challenge is to isolate visual evidence use from generic image damage. Masking any region can disrupt the input, so we compare a human-annotated answer region with an equally sized non-evidence region on the same image. We then restore the clean hidden states at corresponding model layers and token positions during the corrupted run to locate where the additional answer-score loss can be recovered. Across Gemma3, Qwen2.5-VL, and LLaVA1.5 \citep{gemma3,qwen25vl,llava}, answer-region corruption causes greater loss than matched corruption, and this matched-position restoration selectively recovers it. The two interventions establish an evidence-to-answer response at the input and hidden-state levels and give the next challenge a concrete target: the score change produced by removing the annotated evidence.

The second challenge is to determine whether an internal description accounts for the effect it is proposed to explain. A complete hidden state, a fitted sparse reconstruction, selected model channels, and a low-rank subspace are different objects; fitted sparse features are analysis tools rather than components of the base VLM \citep{dunefsky2024transcoders,ameisen2025circuittracing,geiger2024das,makelov2024subspace}. We test each candidate at a specified interface in both directions: inserting its clean value into the masked run should recover answer support, and writing its masked value into the clean run should remove the corresponding support. Each candidate is compared with the appropriate complete intervention at that interface; the full-hidden test uses the natural clean-to-mask score loss itself. The tested compact descriptions do not reproduce their complete references under this criterion. A frozen full-hidden LLaVA interface does: on 60 separately reviewed examples, it restores and removes about 95\% of the mean natural score change. The contrast reveals a \emph{dependence--visibility gap}: visual evidence can affect answer-side computation even when a particular compact lens does not faithfully reproduce its reference effect.

We contribute a region-anchored, bidirectional test of how faithfully an internal description accounts for a visual answer effect. The same annotated region and matched control first establish the effect to explain; paired internal writes then compare each candidate with a complete reference at its own interface. Applied across three VLMs, this test reveals a shared evidence-specific answer response and measures how much of it different internal descriptions reproduce. The independently confirmed LLaVA complete hidden-state interface provides a constructive example of near-complete effect matching at a frozen answer-side location. Together, these results turn the question of whether a circuit is visible into a measurable question about how much of a specified visual effect it explains.
